\begin{definition}[\textit{Степенная функция.}] Пусть $t \in \R$. Рассмотрим на области $G \subset \C \bs[0,+\infty)$ функцию $$w = |z|^te^{it \cdot arg(z)}$$ где $arg(z) \in (0, 2\pi)$. 

Эта функция регулярна на $G$. Причём, при $t \neq 0$ функция однолистна на области $D \subset G$, если $D$ не содержит двух различных точек $z_1, z_2$ таких, что  $z_1 = z_2e^{\frac{2\pi i}{n}}$. 

В частности, при $t > 0$ эта функция осуществляет конформное отображение углов области $G_{0,\phi_0} = \{z : |z| > 0, 0 < arg(z) < \phi_0\}$, где $\phi_0 \leq 2\pi, |t|\phi_0 \leq 2\pi$, на угловую область $G_{0, t\phi_0}$ (первый индекс -- вершина, второй -- величина угла)
\end{definition}

\begin{unnecessary}
\begin{example}
Функция $w = z^2$ отображает
\begin{enumerate}
    \item верхнюю полуплоскость $\{z : Im z > 0\}$ на плоскость с разрезом $\C \bs [0,+\infty)$
    \item полукруг $\{z : |z| < 1, Im z > 0\}$ на круг с разрезом $\{w: |w| < 1\} \bs [0,1)$
    \item полуплоскость $\{z: Im z > a > 0\}$ на внешность параболы $\{w = u + iv: v^2 > 4a^2(u + a^2)\}$
\end{enumerate}
\end{example}
\end{unnecessary}

\begin{definition}[\textit{Экспоненциальная функция.}] Функция $w = e^z$ осуществляет конформное отображение в области $D \subset \C$ тогда и только тогда, когда $D$ не содержит двух различных точек $z_1, z_2$ таких, что $z_1-z_2 = 2n\pi i$
\end{definition}

\begin{unnecessary}
\begin{example}
Функция $w = e^z$ отображает
\begin{enumerate}
    \item полосу $\{z : 0 < Im z < \pi\}$ на верхнюю полуплоскость $\{w : Im w > 0\}$
    \item полуполосу $\{z : Re z < 0, 0 < Im z < \pi\}$ на полукруг $\{w : |w| < 1, Im w > 0\}$
    \item полуполосу $\{z : Re z > 0, 0 < Im z < \pi\}$ на область $\{w : |w| > 1, Im w > 0\}$
\end{enumerate}
\end{example}
\end{unnecessary}

